\subsection{\texorpdfstring{平面 \(E_{s,s'}\) 与由 \(\sigma_{s}\) 和 \(\sigma_{s'}\) 生成的群}{平面 E\_s,s' 与由 sigma\_s 和 sigma\_s' 生成的群}}

在本号中，我们用 \(s\) 与 \(s'\) 表示 \(S\) 的两个满足 \(s \neq s'\) 的元素。置 \(m = m(s, s')\) ，并用 \(\mathrm{E}_{s,s'}\) 表示平面 \(\mathbf{Re}_s \oplus \mathbf{Re}_{s'}\) 。

命题 1. \(B_{M}\) 在 \(E_{s,s'}\) 上的限制是正的，且它非退化当且仅当 \(m\) 有限。

设 \(z = xe_{s} + ye_{s'}\) （其中 \(x, y \in \mathbf{R}\) ）是 \(E_{s,s'}\) 的一个元素。我们有

\[
\begin{array}{r} B _ {M} (z, z) = x ^ {2} - 2 x y \cos \frac {\pi}{m} + y ^ {2} \\ = (x - y. \cos \frac {\pi}{m}) ^ {2} + y ^ {2} \sin^ {2} \frac {\pi}{m}, \end{array}
\]

这表明 \(\mathrm{B}_{\mathsf{M}}\) 在 \(\mathrm{E}_{s,s'}\) 上是正的，并且它在该处非退化当且仅当 \(\sin \frac{\pi}{m} \neq 0\) 。命题得证。

反射 \(\sigma_{s}\) 与 \(\sigma_{s'}\) 保持 \(E_{s,s'}\) 稳定。我们将确定 \(\sigma_{s}\sigma_{s'}\) 在 \(E_{s,s'}\) 上限制的阶。我们区分两种情形：

\begin{enumerate}
\def\labelenumi{\alph{enumi})}
\tightlist
\item
  \(m = +\infty\)
\end{enumerate}

设 \(u = e_s + e_{s'}\) 。我们有 \(\mathrm{B}_{\mathrm{M}}(u, e_s) = \mathrm{B}_{\mathrm{M}}(u, e_{s'}) = 0\) ，故 \(u\) 在 \(\sigma_s\) 与 \(\sigma_{s'}\) 下不变。此外，

\[
\sigma_ {s} \left(\sigma_ {s ^ {\prime}} \left(e _ {s}\right)\right) = \sigma_ {s} \left(e _ {s} + 2 e _ {s ^ {\prime}}\right) = 3 e _ {s} + 2 e _ {s ^ {\prime}} = 2 u + e _ {s},
\]

从而

\[
\left(\sigma_ {s} \sigma_ {s ^ {\prime}}\right) ^ {n} \left(e _ {s}\right) = 2 n u + e _ {s} \text {   for   all   } n \in \mathbf {Z}.
\]

由此可知，\(\sigma_s\sigma_{s'}\) 在 \(\mathrm{E}_{s,s'}\) 上的限制是无限阶的。

\begin{enumerate}
\def\labelenumi{\alph{enumi})}
\setcounter{enumi}{1}
\tightlist
\item
  \(m\) 是有限的。
\end{enumerate}

型 \(B_{M}\) 给 \(E_{s,s'}\) 赋予欧氏平面的结构。由于 \(e_{s}\) 与 \(e_{s'}\) 的标量积等于 \(-\cos\frac{\pi}{m}=\cos\left(\pi-\frac{\pi}{m}\right)\) ，我们可以给 \(E_{s,s'}\) 定向，使得半直线 \(R_{+}e_{s}\) 与 \(R_{+}e_{s'}\) 之间的角等于 \(\pi-\frac{\pi}{m}\) 。若 D 与 \(D'\) 表示正交于 \(e_{s}\) 与 \(e_{s'}\) 的直线，

\[
(\widehat {D ^ {\prime} , D}) = \pi - (\widehat {D , D ^ {\prime}}) = \frac {\pi}{m}.
\]

现在，限制 \(\overline{\sigma}_s\) 与 \(\overline{\sigma}_{s'}\) （即 \(\sigma_s\) 与 \(\sigma_{s'}\) 在 \(E_{s,s'}\) 上的限制）是关于 D 与 \(D'\) 的正交对称。由 §2 第 5 号命题 6 的推论，可知 \(\overline{\sigma}_s\overline{\sigma}_{s'}\) 是角为 \(\frac{2\pi}{m}\) 的旋转。特别地，它的阶是 \(m\) 。

现在回到 E：

命题 2. \(\mathbf{GL}(\mathrm{E})\) 中由 \(\sigma_s\) 与 \(\sigma_{s'}\) 生成的子群是 \(2m(s,s')\) 阶的二面体群。

由于 \(\sigma_s\) 与 \(\sigma_{s'}\) 都是 2 阶的且互不相同，只需证明它们的乘积 \(\sigma_s\sigma_{s'}\) 的阶是 \(m(s,s')\) 。当 \(m(s,s')\) 无限时，这由上面的 a) 得出。当 \(m(s, s')\) 有限时，由命题 1 可知 E 是 \(E_{s,s'}\) 与其正交补 \(V_{s,s'}\) 的直和；由于 \(\sigma_s\) 与 \(\sigma_{s'}\) 在 \(V_{s,s'}\) 上是恒等，且由 b) 知 \(\sigma_s\sigma_{s'}\) 在 \(E_{s,s'}\) 上的限制的阶是 \(m(s, s')\) ，故 \(\sigma_s\sigma_{s'}\) 的阶确实等于 \(m(s, s')\) 。

